%% ma: L12-B18-0013
%% lop: 12
%% chuong: 6
%% bai: 18
%% dang: 12.18.1
%% loai: DS
%% muc: [NB, NB, TH, TH]
%% dap_an: "SSĐĐ"
%% nguon: "(NBV)-10 ĐỀ KIỂM TRA HK2-TOÁN 12-DUNG CHUNG 3 SÁCH.pdf (H01, Đề số 1), Phần 2 Câu 1"
%% kien_thuc: "Bài 18 (xác suất có điều kiện từ số liệu đếm)"
%% trang_thai: chinh-thuc
%% ghi_chu: "Đáp án gốc S S Đ Đ khớp (giải độc lập: P(A) = 0,625; P(B) = 0,4; P(A | B) = 0,75; P(B | A) = 0,48). Trùng ý tưởng với dạng đã duyệt 12.18.1 (xác suất có điều kiện từ số liệu thống kê hai chiều), nhập chính thức. Bỏ ảnh minh họa lớp học (ảnh trang trí)."
\begin{ex}
Lớp $12$A có $40$ học sinh, trong đó có $25$ học sinh tham gia câu lạc bộ Tiếng Anh, $16$ học sinh tham gia câu lạc bộ Toán, $12$ học sinh vừa tham gia câu lạc bộ Tiếng Anh vừa tham gia câu lạc bộ Toán. Chọn ngẫu nhiên $1$ học sinh. Xét các biến cố sau: $A$: ``Học sinh được chọn tham gia câu lạc bộ Tiếng Anh''; $B$: ``Học sinh được chọn tham gia câu lạc bộ Toán''.
\choiceTF
{$P(A)=0{,}4$.}
{$P(B)=0{,}625$.}
{\True $P(A\mid B)=0{,}75$.}
{\True $P(B\mid A)=0{,}48$.}
\loigiai{$P(A)=\dfrac{25}{40}=0{,}625$ nên a) sai. $P(B)=\dfrac{16}{40}=0{,}4$ nên b) sai. $P(AB)=\dfrac{12}{40}=0{,}3$.

$P(A\mid B)=\dfrac{12}{16}=0{,}75$ nên c) đúng. $P(B\mid A)=\dfrac{12}{25}=0{,}48$ nên d) đúng.}
\end{ex}
